\documentclass[12pt]{article}
\usepackage[margin=1in]{geometry}
\usepackage{amsmath,amssymb,amsfonts,amsthm,bm,array,booktabs,enumitem}
\usepackage[colorlinks=true,linkcolor=blue,citecolor=blue,urlcolor=blue]{hyperref}

\numberwithin{equation}{section}

\newtheorem{theorem}{Theorem}
\newtheorem{proposition}{Proposition}
\newtheorem{lemma}{Lemma}
\newtheorem{assumption}{Assumption}
\theoremstyle{definition}
\newtheorem{definition}{Definition}
\newtheorem{remark}{Remark}

\newcommand{\E}{\mathbb E}
\newcommand{\R}{\mathbb R}
\newcommand{\Pp}{\mathbb P}
\newcommand{\ind}{\mathbf 1}
\newcommand{\op}{o_p}
\newcommand{\Op}{O_p}
\newcommand{\N}{\mathcal N}
\newcommand{\LtwoT}{L^2(P_T)}
\newcommand{\LtwoP}{L^2(P)}
\newcommand{\innerT}[2]{\langle #1,#2\rangle_T}
\newcommand{\innerP}[2]{\langle #1,#2\rangle_P}

\begin{document}

\begin{center}
{\Large\bf L2 Geometry Workbook for Profile Sieve EL}\\[6pt]
{\large Three-Part Proof: Geometry, Small Stochastic Bounds, and Merging}\\[6pt]
{\normalsize Solved derivation workbook}\\[12pt]
\end{center}

\tableofcontents
\newpage

\part{Geometric and Convex KKT Body}

\section{Model Setup: Full Geometry in \texorpdfstring{\(L^2(P_T)\)}{L2PT} and \texorpdfstring{\(L^2(P)\)}{L2P}}

The purpose of this workbook is to rewrite the whole proof logic as Hilbert
space geometry. The sample space is not first probability; it is first the
finite empirical Hilbert space
\[
    \LtwoT=(\R^T,\innerT{a}{b}=T^{-1}a'b).
\]
The population limit is the Hilbert space
\[
    \LtwoP,\qquad \innerP{f}{g}=\E\{fg\}.
\]

The model setup is the following chain:
\[
\begin{aligned}
\bm y&=\beta_0\bm x+\bm\mu_0+\bm u
&&\text{model in }\LtwoT,\\
\mu_{0,t}&=m_{j_0(t)}(W_{t-1}),\quad
j_0(t)=j\iff k_{j,0}<t\le k_{j+1,0}
&&\text{true nuisance vector},\\
\N_{K,\Lambda,T}&=\operatorname{col}(\bm Q_{K,\Lambda})
&&\text{nuisance subspace},\\
\bm M_{K,\Lambda}&=I-\Pi_{K,\Lambda,T}
&&\text{residual-maker},\\
S_T(\beta,\Lambda)&=\sqrt T\innerT{\bm M_{K,\Lambda}(\bm y-\beta\bm x)}
{\bm M_{K,\Lambda}\bm w}
&&\text{projected score}.
\end{aligned}
\]
Here \(\bm\mu_0\) is exactly the stacked nonparametric nuisance part of the
model. If there are no nuisance breaks, \(\mu_{0,t}=m(W_{t-1})\). With breaks,
\(\mu_{0,t}=m_j(W_{t-1})\) on regime \(j\), so \(\bm\mu_0\) is not one stable
function evaluated over the whole sample; it is the true piecewise nuisance
surface evaluated observation by observation.

EL then uses the scalar products
\[
    Z_t(\beta,\Lambda)
    =\{\bm M_{K,\Lambda}(\bm y-\beta\bm x)\}_t\{\bm M_{K,\Lambda}\bm w\}_t
\]
to form a one-dimensional empirical likelihood. EL is not the Hilbert
projection; EL is the likelihood wrapper around the projected score.

\begin{center}
\fbox{\parbox{0.92\textwidth}{
\textbf{Three-part proof map.} Part I is finite-sample geometry: profiling,
residual scores, and EL are all KKT/convex objects in \(\LtwoT\). Part II is
small stochastic verification: Gram stability, sieve and break errors, CLT,
variance, max-score, and convex-hull size. Part III merges them: consistency
comes from the moment bridge plus unique population zero, while Wilks comes
from the EL KKT quadratic reduction plus the standardized score CLT.}}
\end{center}

\section{Finite-Sample Hilbert Space: Viewing the Entire Model in \texorpdfstring{\(L^2(P_T)\)}{L2PT}}

\subsection{The empirical space}

Let
\[
    \Omega_T=\{1,\ldots,T\},
    \qquad
    P_T=T^{-1}\sum_{t=1}^T\delta_t.
\]
Every sample vector \(\bm a=(a_1,\ldots,a_T)'\) is a function
\(a_T:\Omega_T\to\R\), with \(a_T(t)=a_t\). The empirical inner product is
\[
    \innerT{\bm a}{\bm b}
    =\int a_T b_T\,dP_T
    =T^{-1}\sum_{t=1}^Ta_tb_t
    =T^{-1}\bm a'\bm b.
\]
Thus all variables in the sample are elements of \(\LtwoT\):
\[
    \bm y,\bm x,\bm w,\bm u,\bm\mu_0\in\LtwoT.
\]

\subsection{The full finite-sample model}

The break-aware model is the vector identity in \(\LtwoT\)
\begin{equation}
    \bm y=\beta_0\bm x+\bm\mu_0+\bm u,
    \label{eq:l2pt-model}
\end{equation}
where
\[
    \mu_{0,t}=\sum_{j=0}^{q_0}m_j(W_{t-1})\ind\{k_{j,0}<t\le k_{j+1,0}\}.
\]
So the full model says: \(\bm y\) lies in the affine set
\[
    \beta_0\bm x+\bm\mu_0+\bm u.
\]
The inferential problem is to remove \(\bm\mu_0\), not by estimating it as a
main parameter, but by projecting away its nuisance subspace.

\subsection{Candidate nuisance subspaces}

For a candidate partition \(\Lambda=(k_1,\ldots,k_q)\), define the block sieve
matrix
\[
    \bm Q_{K,\Lambda}
    =[
      \bm D_0(\Lambda)\bm P,
      \ldots,
      \bm D_q(\Lambda)\bm P
    ].
\]
The sample nuisance subspace is
\begin{equation}
    \N_{K,\Lambda,T}
    =\operatorname{col}(\bm Q_{K,\Lambda})\subset\LtwoT.
    \label{eq:sample-nuisance-space}
\end{equation}
An element of this space has the form
\[
    g_T(t)=\sum_{j=0}^q\ind\{t\in I_j(\Lambda)\}\bm P_K(W_{t-1})'c_j.
\]
This is the finite-sample version of a regime-specific nonparametric nuisance
function.

Let \(j_\Lambda(t)\) denote the regime index induced by \(\Lambda\). For a
regime-wise nuisance tuple \(h=(h_0,\ldots,h_q)\), define its stacked sample
vector by
\[
    \bm h_{\Lambda,T}
    =\big(h_{j_\Lambda(1)}(W_0),\ldots,
          h_{j_\Lambda(T)}(W_{T-1})\big)'.
\]
This vector is the empirical version of the candidate nuisance direction.

\begin{lemma}[Uniform finite-sieve approximation in \texorpdfstring{\(L^2(P_T)\)}{L2PT}]
\label{lem:sieve-decomposition}
Fix \((K,\Lambda,T)\) and a nuisance class \(\mathcal H\). For each
\(h\in\mathcal H\), let \(\bm h_{\Lambda,T}\in\LtwoT\) be the stacked nuisance
vector defined above. Define the finite-sieve approximation radius
\begin{equation}
    \rho_{K,T}(\mathcal H,\Lambda)
    =\sup_{h\in\mathcal H}
      \inf_{\gamma\in\R^{(q+1)K}}
      \|\bm h_{\Lambda,T}-\bm Q_{K,\Lambda}\gamma\|_T.
    \label{eq:rho-sieve-radius}
\end{equation}
If \(\rho_{K,T}(\mathcal H,\Lambda)=\op(1)\), then the desired uniform
\(\LtwoT\) approximation holds:
\begin{equation}
    \sup_{h\in\mathcal H}
    \|\bm h_{\Lambda,T}-\bm Q_{K,\Lambda}\gamma_{K,\Lambda}(h)\|_T
    =\sup_{h\in\mathcal H}\|\bm r_{K,\Lambda,T}(h)\|_T
    =\rho_{K,T}(\mathcal H,\Lambda)
    =\op(1).
    \label{eq:uniform-sieve-approximation}
\end{equation}
Here, for each \(h\),
\[
    \gamma_{K,\Lambda}(h)
    \in\operatorname*{argmin}_{\gamma\in\R^{(q+1)K}}
    \|\bm h_{\Lambda,T}-\bm Q_{K,\Lambda}\gamma\|_T^2,
\]
\[
    \bm h_{K,\Lambda,T}
    =\bm Q_{K,\Lambda}\gamma_{K,\Lambda}(h),
    \qquad
    \bm r_{K,\Lambda,T}(h)
    =\bm h_{\Lambda,T}-\bm h_{K,\Lambda,T}.
\]
The ambient space is \(\LtwoT\cong\R^T\), so \(\bm h_{\Lambda,T}\),
\(\bm h_{K,\Lambda,T}\), and \(\bm r_{K,\Lambda,T}(h)\) are all \(T\times1\)
vectors. The design matrix and coefficient vector have dimensions
\[
    \bm Q_{K,\Lambda}\in\R^{T\times(q+1)K},
    \qquad
    \gamma_{K,\Lambda}(h)\in\R^{(q+1)K}.
\]
The finite nuisance space has dimension
\[
    d_{K,\Lambda,T}=\operatorname{rank}(\bm Q_{K,\Lambda})\le(q+1)K,
\]
and
\begin{equation}
    \bm h_{\Lambda,T}
    =\bm Q_{K,\Lambda}\gamma_{K,\Lambda}(h)
     +\bm r_{K,\Lambda,T}(h),
    \label{eq:sieve-decomposition}
\end{equation}
where
\[
    \bm Q_{K,\Lambda}\gamma_{K,\Lambda}(h)\in\N_{K,\Lambda,T}\subset\LtwoT,
    \qquad
    \bm r_{K,\Lambda,T}(h)\in\N_{K,\Lambda,T}^{\perp_T}\subset\LtwoT.
\]
Thus \(\bm r_{K,\Lambda,T}(h)\) is a \(T\times1\) approximation-residual vector
living in a \((T-d_{K,\Lambda,T})\)-dimensional orthogonal complement, and
\[
    \bm Q_{K,\Lambda}'\bm r_{K,\Lambda,T}(h)=0.
\]
\end{lemma}

\begin{proof}
The column space \(\N_{K,\Lambda,T}=\operatorname{col}(\bm Q_{K,\Lambda})\) is a
finite-dimensional closed subspace of \(\LtwoT\). Therefore the Hilbert
projection theorem gives the closest finite-sieve approximation to every
\(\bm h_{\Lambda,T}\):
\[
    \Pi_{K,\Lambda,T}\bm h_{\Lambda,T}
    \in\N_{K,\Lambda,T}.
\]
Because every element of \(\N_{K,\Lambda,T}\) is \(\bm Q_{K,\Lambda}\gamma\), the
projection can be written as
\[
    \Pi_{K,\Lambda,T}\bm h_{\Lambda,T}
    =\bm Q_{K,\Lambda}\gamma_{K,\Lambda}(h)
\]
for any least-squares minimizer \(\gamma_{K,\Lambda}(h)\). If
\(\bm Q_{K,\Lambda}\) is rank deficient, the coefficient vector need not be
unique, but the fitted vector is unique.

Define
\[
    \bm r_{K,\Lambda,T}(h)
    =\bm h_{\Lambda,T}-\Pi_{K,\Lambda,T}\bm h_{\Lambda,T}.
\]
The projection theorem gives
\[
    \innerT{\bm r_{K,\Lambda,T}(h)}{\bm v}=0
    \qquad\text{for every }\bm v\in\N_{K,\Lambda,T},
\]
which implies \(\bm Q_{K,\Lambda}'\bm r_{K,\Lambda,T}(h)=0\). This proves the
orthogonal decomposition \eqref{eq:sieve-decomposition} and the location
\(\bm r_{K,\Lambda,T}(h)\in\N_{K,\Lambda,T}^{\perp_T}\).

By the definition of orthogonal projection,
\[
    \|\bm r_{K,\Lambda,T}(h)\|_T
    =\inf_{\gamma\in\R^{(q+1)K}}
      \|\bm h_{\Lambda,T}-\bm Q_{K,\Lambda}\gamma\|_T.
\]
Taking the supremum over \(h\in\mathcal H\) gives exactly
\eqref{eq:uniform-sieve-approximation}. Hence the mathematical task behind this
lemma is to verify \(\rho_{K,T}(\mathcal H,\Lambda)=\op(1)\) from the chosen
basis, smoothness class, and growth rate of \(K\). An arbitrary nuisance
direction is not assumed to be exactly spanned by the sieve basis; it is
approximated by its finite-sieve projection, and the leftover is
\(\bm r_{K,\Lambda,T}(h)\).
\end{proof}
\subsection{Orthogonal projection}

Let \(\Pi_{K,\Lambda,T}\) be the \(\LtwoT\)-orthogonal projection onto
\(\N_{K,\Lambda,T}\). In coordinates,
\[
    \Pi_{K,\Lambda,T}
    =\bm Q_{K,\Lambda}(\bm Q_{K,\Lambda}'\bm Q_{K,\Lambda})^\dagger
      \bm Q_{K,\Lambda}'.
\]
The residual-maker is
\begin{equation}
    \bm M_{K,\Lambda}=I-\Pi_{K,\Lambda,T}.
    \label{eq:sample-residual-maker}
\end{equation}
It satisfies the projection identities
\[
    \bm M_{K,\Lambda}'=\bm M_{K,\Lambda},
    \qquad
    \bm M_{K,\Lambda}^2=\bm M_{K,\Lambda},
    \qquad
    \bm Q_{K,\Lambda}'\bm M_{K,\Lambda}=0.
\]

\begin{lemma}[Projection identities for the residual-maker]
\label{lem:projection-identities}
Let \(Q=\bm Q_{K,\Lambda}\), \(G=Q'Q\),
\[
    \Pi=QG^\dagger Q',
    \qquad
    M=I-\Pi,
\]
where \(G^\dagger\) is the Moore-Penrose inverse. Then \(\Pi\) and \(M\) are
symmetric and idempotent, and \(M\) annihilates the sieve nuisance space:
\[
    M'=M,
    \qquad
    M^2=M,
    \qquad
    MQ=0,
    \qquad
    Q'M=0.
\]
\end{lemma}

\begin{proof}
Since \(G=Q'Q\) is symmetric, its Moore-Penrose inverse \(G^\dagger\) is also
symmetric. Hence
\[
    \Pi'=\{QG^\dagger Q'\}'=QG^\dagger Q'=\Pi.
\]
For idempotence, use the Moore-Penrose identity
\(G^\dagger GG^\dagger=G^\dagger\):
\[
    \Pi^2
    =QG^\dagger Q'QG^\dagger Q'
    =QG^\dagger GG^\dagger Q'
    =QG^\dagger Q'
    =\Pi.
\]
Therefore
\[
    M'=I-\Pi'=M,
    \qquad
    M^2=(I-\Pi)^2=I-2\Pi+\Pi^2=I-\Pi=M.
\]
It remains to show that \(M\) kills the columns of \(Q\). Since
\(\mathcal R(G)=\mathcal R(Q')\) and \(G^\dagger G\) is the orthogonal projector
onto \(\mathcal R(G)\),
\[
    QG^\dagger G=Q.
\]
Thus
\[
    \Pi Q=QG^\dagger Q'Q=QG^\dagger G=Q,
    \qquad
    MQ=(I-\Pi)Q=0.
\]
Finally, because \(M\) is symmetric, \(Q'M=(MQ)'=0\). This proves all projection
identities.
\end{proof}

\section{Profiling Is Orthogonal Projection}

For a trial value \(\beta\), define the partial outcome
\[
    \bm y_\beta^\circ=\bm y-\beta\bm x.
\]
It is the original outcome \(\bm y\) after subtracting the trial linear component
\(\beta\bm x\).

Profiling the nuisance means finding the closest nuisance vector to the partial
outcome \(\bm y_\beta^\circ\) in \(\N_{K,\Lambda,T}\):
\[
    \widehat{\bm\mu}_{\beta,\Lambda}
    =\arg\min_{\bm g\in\N_{K,\Lambda,T}}\|\bm y_\beta^\circ-\bm g\|_T^2.
\]
By Hilbert projection,
\[
    \widehat{\bm\mu}_{\beta,\Lambda}=\Pi_{K,\Lambda,T}\bm y_\beta^\circ,
    \qquad
    \widehat{\bm u}(\beta,\Lambda)=\bm y_\beta^\circ-\widehat{\bm\mu}_{\beta,\Lambda}
    =\bm M_{K,\Lambda}\bm y_\beta^\circ.
\]

\begin{lemma}[Profile normal equation]
For every \(\bm h\in\N_{K,\Lambda,T}\),
\[
    \innerT{\widehat{\bm u}(\beta,\Lambda)}{\bm h}=0.
\]
Equivalently,
\[
    \bm Q_{K,\Lambda}'\widehat{\bm u}(\beta,\Lambda)=0.
\]
\end{lemma}

\begin{proof}
This is the KKT condition for a least-squares projection in \(\LtwoT\). If
\(\bm h=\bm Q_{K,\Lambda}\bm a\), then
\[
    \innerT{\widehat{\bm u}}{\bm h}
    =T^{-1}\widehat{\bm u}'\bm Q_{K,\Lambda}\bm a
    =T^{-1}(\bm Q_{K,\Lambda}'\widehat{\bm u})'\bm a=0.
\]
\end{proof}

\subsection{Break search as subspace selection}

For fixed \(\beta\), the partial outcome is
\[
    \bm y_\beta^\circ=\bm y-\beta\bm x.
\]
For a candidate partition \(\Lambda=(k_1,\ldots,k_q)\), set \(k_0=0\),
\(k_{q+1}=T\), and
\[
    I_j(\Lambda)=\{t:k_j<t\le k_{j+1}\}.
\]
The block sieve matrix is
\[
    \bm Q_{K,\Lambda}
    =[\bm D_0(\Lambda)\bm P,\ldots,\bm D_q(\Lambda)\bm P].
\]
For \(c=(c_0',\ldots,c_q')'\), the fitted nuisance vector is
\(\bm Q_{K,\Lambda}c\), with component
\[
    (\bm Q_{K,\Lambda}c)_t
    =\bm P_K(W_{t-1})'c_j
    \qquad\text{if }t\in I_j(\Lambda).
\]
Thus minimizing over \(c\) is exactly the same as choosing a separate sieve
coefficient vector \(c_j\) on each segment.

\begin{lemma}[Profile RSS is a Hilbert distance]
\label{lem:rss-hilbert-distance}
For fixed \((\beta,\Lambda)\), define the segment cost
\[
    C_K(s,e;\beta)
    =\min_{c\in\R^K}
      \sum_{t=s}^e
      \{y_t-\beta x_t-\bm P_K(W_{t-1})'c\}^2,
\]
and the profile RSS
\[
    RSS_K(\beta,\Lambda)
    =\sum_{j=0}^q C_K(k_j+1,k_{j+1};\beta).
\]
Then
\begin{equation}
\begin{aligned}
    RSS_K(\beta,\Lambda)
    &=\min_{c\in\R^{(q+1)K}}
      \|\bm y_\beta^\circ-\bm Q_{K,\Lambda}c\|_2^2 \\
    &=\inf_{\bm g\in\N_{K,\Lambda,T}}
      \|\bm y_\beta^\circ-\bm g\|_2^2 \\
    &=T\operatorname{dist}_T^2
      (\bm y_\beta^\circ,\N_{K,\Lambda,T}) \\
    &=T\|\bm M_{K,\Lambda}\bm y_\beta^\circ\|_T^2.
\end{aligned}
\label{eq:rss-distance-equivalence}
\end{equation}
\end{lemma}

\begin{proof}
Because each segment \(I_j(\Lambda)\) has its own coefficient vector \(c_j\),
\[
\begin{aligned}
    RSS_K(\beta,\Lambda)
    &=\min_{c_0,\ldots,c_q}
      \sum_{j=0}^q\sum_{t\in I_j(\Lambda)}
      \{y_t-\beta x_t-\bm P_K(W_{t-1})'c_j\}^2 \\
    &=\min_{c\in\R^{(q+1)K}}
      \|\bm y-\beta\bm x-\bm Q_{K,\Lambda}c\|_2^2 \\
    &=\min_{c\in\R^{(q+1)K}}
      \|\bm y_\beta^\circ-\bm Q_{K,\Lambda}c\|_2^2.
\end{aligned}
\]
Since \(\N_{K,\Lambda,T}=\operatorname{col}(\bm Q_{K,\Lambda})\), every
\(\bm g\in\N_{K,\Lambda,T}\) can be written as
\(\bm g=\bm Q_{K,\Lambda}c\). Therefore
\[
    \min_c\|\bm y_\beta^\circ-\bm Q_{K,\Lambda}c\|_2^2
    =\inf_{\bm g\in\N_{K,\Lambda,T}}
      \|\bm y_\beta^\circ-\bm g\|_2^2.
\]
The Euclidean norm and the empirical Hilbert norm satisfy
\(\|\bm a\|_2^2=T\|\bm a\|_T^2\). Hence
\[
    \inf_{\bm g\in\N_{K,\Lambda,T}}
      \|\bm y_\beta^\circ-\bm g\|_2^2
    =T\operatorname{dist}_T^2(\bm y_\beta^\circ,\N_{K,\Lambda,T}).
\]
Finally, \(\bm M_{K,\Lambda}\) is the residual-maker for the orthogonal
projection onto \(\N_{K,\Lambda,T}\), so
\[
    \operatorname{dist}_T^2(\bm y_\beta^\circ,\N_{K,\Lambda,T})
    =\|\bm M_{K,\Lambda}\bm y_\beta^\circ\|_T^2.
\]
This proves \eqref{eq:rss-distance-equivalence}.
\end{proof}

Consequently, for known order \(q\), the break estimator is
\[
\begin{aligned}
    \widehat\Lambda_q(\beta)
    &=\arg\min_{\Lambda\in\mathcal L_q(\epsilon)} RSS_K(\beta,\Lambda) \\
    &=\arg\min_{\Lambda\in\mathcal L_q(\epsilon)}
      \operatorname{dist}_T^2(\bm y_\beta^\circ,\N_{K,\Lambda,T}).
\end{aligned}
\]
Unknown order adds a penalty to the same distance:
\[
    \widehat q(\beta)
    =\arg\min_{0\le q\le\bar q}
     \left\{
       T\operatorname{dist}_T^2
       (\bm y_\beta^\circ,\N_{K,\widehat\Lambda_q(\beta),T})
       +\operatorname{pen}_T(q,K)
     \right\}.
\]
Thus break search is subspace selection: each candidate \(\Lambda\) creates a
candidate nuisance subspace \(\N_{K,\Lambda,T}\), and RSS picks the subspace
closest to the partial outcome \(\bm y_\beta^\circ\).
\section{Residualized Score Direction}

The raw score weight \(\bm w\in\LtwoT\) may contain nuisance directions. The
geometric score direction is its nuisance-orthogonal component
\[
    \bm v_\Lambda=\bm w^c(\Lambda)=\bm M_{K,\Lambda}\bm w.
\]
Then
\[
    \bm v_\Lambda\in\N_{K,\Lambda,T}^{\perp_T},
    \qquad
    \innerT{\bm v_\Lambda}{\bm h}=0
    \quad \forall \bm h\in\N_{K,\Lambda,T}.
\]

The sample profiled moment is the Hilbert inner product
\begin{equation}
    \Psi_T(\beta,\Lambda)
    =\innerT{\bm M_{K,\Lambda}(\bm y-\beta\bm x)}
      {\bm M_{K,\Lambda}\bm w}.
    \label{eq:sample-profile-moment}
\end{equation}
The scaled score is
\[
    S_T(\beta,\Lambda)=\sqrt T\,\Psi_T(\beta,\Lambda).
\]

\begin{lemma}[Finite-sample nuisance cancellation]
If \(\bm g\in\N_{K,\Lambda,T}\), then
\[
    \innerT{\bm M_{K,\Lambda}\bm g}{\bm M_{K,\Lambda}\bm w}=0.
\]
\end{lemma}

\begin{proof}
Because \(\bm g\in\N_{K,\Lambda,T}\), \(\bm M_{K,\Lambda}\bm g=0\). Hence the
inner product is zero. Equivalently,
\((\bm Qc)'\bm M_{K,\Lambda}\bm w=0\) for every \(c\).
\end{proof}

\section{Known True Partition: Wilks Logic in Geometry}

Set \(\Lambda=\Lambda_0\), and abbreviate
\[
    \bm Q_0=\bm Q_{K,\Lambda_0},
    \qquad
    \bm M_0=\bm M_{K,\Lambda_0},
    \qquad
    \bm v_0=\bm M_0\bm w.
\]
Decompose the true nuisance into its finite-sieve projection and approximation
residual:
\[
    \bm\mu_0=\bm Q_0\bm c_0+\bm r_0,
    \qquad
    \bm Q_0\bm c_0\in\N_{K,\Lambda_0,T},
    \qquad
    \bm r_0\in\N_{K,\Lambda_0,T}^{\perp_T}.
\]

\begin{lemma}[True-partition score decomposition]
\label{lem:true-partition-score-decomposition}
At \(\beta=\beta_0\) and \(\Lambda=\Lambda_0\), the projected score satisfies
\begin{equation}
    S_T(\beta_0,\Lambda_0)
    =T^{-1/2}\bm u'\bm v_0
     +T^{-1/2}\bm r_0'\bm v_0.
    \label{eq:true-partition-score-decomposition}
\end{equation}
The first term is the stochastic score. The second term is the sieve
approximation-bias term, and this is the term controlled by approximation
rates:
\begin{equation}
    T^{-1/2}\bm r_0'\bm v_0
    =T^{-1/2}\bm r_0'\bm M_0\bm w
    =\op(1).
    \label{eq:true-partition-approximation-bias}
\end{equation}
\end{lemma}

\begin{proof}
At the true parameter, the partial outcome is
\[
    \bm y_{\beta_0}^\circ
    =\bm y-\beta_0\bm x
    =\bm\mu_0+\bm u
    =\bm Q_0\bm c_0+\bm r_0+\bm u.
\]
Apply the true-partition residual-maker \(\bm M_0\). Since \(\bm Q_0\bm c_0\) is
in \(\N_{K,\Lambda_0,T}\), Lemma~\ref{lem:projection-identities} gives
\(\bm M_0\bm Q_0\bm c_0=0\). Therefore
\[
    \bm M_0\bm y_{\beta_0}^\circ
    =\bm M_0\bm u+\bm M_0\bm r_0.
\]
The projected score is
\[
\begin{aligned}
    S_T(\beta_0,\Lambda_0)
    &=\sqrt T\innerT{\bm M_0\bm y_{\beta_0}^\circ}{\bm M_0\bm w} \\
    &=T^{-1/2}(\bm M_0\bm u+\bm M_0\bm r_0)'\bm v_0.
\end{aligned}
\]
Because \(\bm M_0\) is symmetric and idempotent and \(\bm v_0=\bm M_0\bm w\),
we have \(\bm M_0\bm v_0=\bm v_0\). Hence
\[
    (\bm M_0\bm u)'\bm v_0=\bm u'\bm v_0,
    \qquad
    (\bm M_0\bm r_0)'\bm v_0=\bm r_0'\bm v_0.
\]
Substituting these two identities gives
\eqref{eq:true-partition-score-decomposition}. The term
\(T^{-1/2}\bm u'\bm v_0\) is handled by the score CLT and variance consistency.
The term \(T^{-1/2}\bm r_0'\bm v_0\) is the only approximation-rate term in this
decomposition, because it is the only term containing the sieve residual
\(\bm r_0\).
\end{proof}
\section{Omitted Breaks: False Predictivity as a Projection Drift}

Let \(\N_{K,0,T}\) be the stable/no-break nuisance space and let
\(\bm M_{\mathrm{st}}\) be the residual-maker for orthogonal projection onto
\(\N_{K,0,T}\). Define
\[
    \bm v_{\mathrm{st}}=\bm M_{\mathrm{st}}\bm w,
    \qquad
    S_T^{\mathrm{st}}(\beta)
    =\sqrt T\innerT{\bm M_{\mathrm{st}}(\bm y-\beta\bm x)}
                         {\bm M_{\mathrm{st}}\bm w}.
\]

\begin{proposition}[Omitted-break projection drift]
\label{prop:omitted-break-projection-drift}
At the true parameter \(\beta_0\),
\begin{equation}
    S_T^{\mathrm{st}}(\beta_0)
    =
    T^{-1/2}\bm u'\bm v_{\mathrm{st}}
    +B_T,
    \qquad
    B_T=T^{-1/2}\bm\mu_0'\bm v_{\mathrm{st}}.
    \label{eq:stable-score-drift-decomposition}
\end{equation}
Equivalently,
\begin{equation}
    B_T
    =
    \sqrt T
    \innerT{\bm M_{\mathrm{st}}\bm\mu_0}
           {\bm M_{\mathrm{st}}\bm w}.
    \label{eq:stable-score-drift-inner-product}
\end{equation}
If
\[
    \innerT{\bm M_{\mathrm{st}}\bm\mu_0}
           {\bm M_{\mathrm{st}}\bm w}
    \to b\ne0,
\]
then \(B_T/\sqrt T\to b\), so the stable score has deterministic drift of
order \(\sqrt T\) at the true parameter.
\end{proposition}

\begin{proof}
At \(\beta_0\),
\[
    \bm y-\beta_0\bm x=\bm\mu_0+\bm u.
\]
Applying the stable residual-maker gives
\[
    \bm M_{\mathrm{st}}(\bm y-\beta_0\bm x)
    =
    \bm M_{\mathrm{st}}\bm u
    +
    \bm M_{\mathrm{st}}\bm\mu_0.
\]
Therefore
\[
\begin{aligned}
    S_T^{\mathrm{st}}(\beta_0)
    &=
    T^{-1/2}
    (\bm M_{\mathrm{st}}\bm u+\bm M_{\mathrm{st}}\bm\mu_0)'
    \bm v_{\mathrm{st}}  \\
    &=
    T^{-1/2}\bm u'\bm v_{\mathrm{st}}
    +
    T^{-1/2}\bm\mu_0'\bm v_{\mathrm{st}},
\end{aligned}
\]
because \(\bm M_{\mathrm{st}}\) is symmetric and idempotent and
\(\bm v_{\mathrm{st}}=\bm M_{\mathrm{st}}\bm w\). This proves
\eqref{eq:stable-score-drift-decomposition}. The same symmetry and idempotence
give
\[
    T^{-1/2}\bm\mu_0'\bm v_{\mathrm{st}}
    =
    T^{-1/2}\bm\mu_0'\bm M_{\mathrm{st}}\bm w
    =
    \sqrt T
    \innerT{\bm M_{\mathrm{st}}\bm\mu_0}
           {\bm M_{\mathrm{st}}\bm w},
\]
which proves \eqref{eq:stable-score-drift-inner-product}. If the empirical
inner product converges to \(b\ne0\), the drift is \(B_T=\sqrt T\,b+o(\sqrt T)\).
\end{proof}

As a toy example, suppose there are no covariates and the stable nuisance space
contains only constants:
\[
    \N_{K,0,T}=\operatorname{span}\{\bm 1_T\}.
\]
Let the true nuisance contain one level shift,
\[
    \mu_{0,t}=\delta\ind\{t>k_0\},
    \qquad
    \pi_T=T^{-1}\sum_{t=1}^T\ind\{t>k_0\}.
\]
The stable projection removes only the global mean. Hence
\[
    (\bm M_{\mathrm{st}}\bm\mu_0)_t
    =
    \delta\{\ind(t>k_0)-\pi_T\}.
\]
If the residualized score direction is aligned with this step, for instance
\(\bm M_{\mathrm{st}}\bm w=\bm M_{\mathrm{st}}\bm\mu_0\), then
\[
    \innerT{\bm M_{\mathrm{st}}\bm\mu_0}
           {\bm M_{\mathrm{st}}\bm w}
    =
    \|\bm M_{\mathrm{st}}\bm\mu_0\|_T^2
    =
    \delta^2\pi_T(1-\pi_T).
\]
When \(\delta\ne0\) and \(\pi_T\to\pi\in(0,1)\), the limit is
\(\delta^2\pi(1-\pi)\ne0\). The stable score therefore reports a false signal
even at \(\beta_0\).

The correct break-aware block sieve changes the projection target. Under the
true partition,
\[
    \bm\mu_0=\bm Q_0\bm c_0+\bm r_0,
    \qquad
    \bm Q_0\bm c_0\in\N_{K,\Lambda_0,T},
    \qquad
    \|\bm r_0\|_T=o_p(1),
\]
so \(\bm M_0\bm\mu_0=\bm r_0\) is small. Under the stable space,
\(\bm M_{\mathrm{st}}\bm\mu_0\) need not be small. The block-sieve projection
removes the low-frequency break component before testing \(\beta\); otherwise
omitted nuisance breaks can enter the score as false predictivity.
\section{Estimated Partition}

The estimated break procedure chooses a random subspace
\[
    \widehat\N_T=\N_{K,\widehat\Lambda,T},
    \qquad
    \widehat M_T=\bm M_{K,\widehat\Lambda}.
\]
The feasible score is
\[
    \widehat S_T(\beta_0)
    =\sqrt T\innerT{\widehat M_T(\bm y-\beta_0\bm x)}{\widehat M_T\bm w}.
\]
The true-partition score is
\[
    S_T^0(\beta_0)
    =\sqrt T\innerT{\bm M_0(\bm y-\beta_0\bm x)}{\bm M_0\bm w},
    \qquad
    \bm M_0=\bm M_{K,\Lambda_0}.
\]
Because each residual-maker is symmetric and idempotent,
\[
    \widehat S_T(\beta_0)=T^{-1/2}(\bm\mu_0+\bm u)'\widehat M_T\bm w,
    \qquad
    S_T^0(\beta_0)=T^{-1/2}(\bm\mu_0+\bm u)'\bm M_0\bm w.
\]
Therefore the exact estimated-partition score error is
\begin{equation}
\begin{aligned}
    \widehat S_T(\beta_0)-S_T^0(\beta_0)
    &=T^{-1/2}\bm\mu_0'(\widehat M_T-\bm M_0)\bm w \\
    &\quad+T^{-1/2}\bm u'(\widehat M_T-\bm M_0)\bm w.
\end{aligned}
\label{eq:estimated-partition-score-decomposition}
\end{equation}
The first term is local sieve-bias perturbation. The second term is local noise
perturbation. Thus the estimated-partition Wilks argument needs
\begin{equation}
    T^{-1/2}\bm\mu_0'(\widehat M_T-\bm M_0)\bm w=\op(1),
    \qquad
    T^{-1/2}\bm u'(\widehat M_T-\bm M_0)\bm w=\op(1).
    \label{eq:estimated-partition-score-pieces}
\end{equation}
For the variance, the corresponding requirement is
\begin{equation}
    \widehat V_T-V_T^0=\op(1),
    \qquad
    \widehat V_T=P_Tg_t(\beta_0,\widehat\Lambda)^2,
    \qquad
    V_T^0=P_Tg_t(\beta_0,\Lambda_0)^2.
    \label{eq:estimated-partition-variance-replacement}
\end{equation}
Equations \eqref{eq:estimated-partition-score-decomposition}--\eqref{eq:estimated-partition-variance-replacement}
are the exact bridge from estimated breaks to the true-partition geometry.

The reason this is weaker than operator-norm convergence is visible from
\eqref{eq:estimated-partition-score-decomposition}. We do not need
\(\|\widehat M_T-\bm M_0\|_{op}\to0\). We only need the difference to be small
when applied to the score direction \(\bm w\) and paired with \(\bm\mu_0\) and
\(\bm u\), plus the analogous second-moment replacement.

The localization proof uses the RSS distance identity from
Lemma~\ref{lem:rss-hilbert-distance}. At \(\beta_0\), set
\[
    \Delta_\Lambda=\bm M_{K,\Lambda}-\bm M_0,
    \qquad
    \bm y_{\beta_0}^\circ=\bm\mu_0+\bm u.
\]
Then
\begin{equation}
\begin{aligned}
    RSS_K(\beta_0,\Lambda)-RSS_K(\beta_0,\Lambda_0)
    &=(\bm y_{\beta_0}^\circ)'\Delta_\Lambda\bm y_{\beta_0}^\circ \\
    &=\bm\mu_0'\Delta_\Lambda\bm\mu_0
      +2\bm\mu_0'\Delta_\Lambda\bm u
      +\bm u'\Delta_\Lambda\bm u.
\end{aligned}
\label{eq:rss-localization-decomposition}
\end{equation}
The deterministic term is the geometric drift from choosing the wrong nuisance
subspace. The two stochastic terms must be uniformly smaller than that drift on
wrong partitions. Part II states the primitive concentration conditions and then
uses \eqref{eq:rss-localization-decomposition} to obtain break localization and
order selection.
\section{EL as a Finite-Sample Convex KKT Problem}

For fixed \((K,\Lambda,\beta)\), everything in this section is finite-sample
algebra in \(\LtwoT\). Define the profiled residual, score direction, and scalar
moment array
\[
    \bm r_\beta=\bm M_{K,\Lambda}(\bm y-\beta\bm x),
    \qquad
    \bm a=\bm M_{K,\Lambda}\bm w,
    \qquad
    g_t(\beta,\Lambda)=r_{\beta,t}a_t.
\]
Then
\[
    \Psi_T(\beta,\Lambda)=P_Tg(\beta,\Lambda)
    =T^{-1}\sum_{t=1}^Tg_t(\beta,\Lambda)
    =\innerT{\bm r_\beta}{\bm a}.
\]
Thus the sieve basis \(Q_{K,\Lambda}\) has already represented the nuisance
feasible directions: \(\operatorname{col}(Q_{K,\Lambda})\) is the tangent plane
for nuisance variation, and \(M_{K,\Lambda}\) keeps only the orthogonal score
component.

EL adds one more finite-dimensional convex problem. The primal feasible set is
the probability simplex cut by the affine moment hyperplane
\[
    \mathcal C_T(\beta,\Lambda)
    =\left\{\pi\in\R^T_+:
       \sum_{t=1}^T\pi_t=1,
       \quad \sum_{t=1}^T\pi_tg_t(\beta,\Lambda)=0
     \right\}.
\]
The empirical likelihood problem is equivalently
\begin{equation}
    \min_{\pi\in\mathcal C_T(\beta,\Lambda)}
    -\sum_{t=1}^T\log(T\pi_t).
    \label{eq:el-primal-convex}
\end{equation}
This is a convex barrier objective over a convex feasible set. Its Lagrangian is
\[
    \mathcal L(\pi,\eta,\lambda)
    =-\sum_{t=1}^T\log(T\pi_t)
      +\eta\left(\sum_{t=1}^T\pi_t-1\right)
      +T\lambda\sum_{t=1}^T\pi_tg_t.
\]
The KKT first-order condition gives
\[
    -\frac{1}{\pi_t}+\eta+T\lambda g_t=0.
\]
After normalizing by \(\sum_t\pi_t=1\), the EL weights are
\begin{equation}
    \widehat\pi_t(\lambda)
    =\frac{1}{T\{1+\lambda g_t(\beta,\Lambda)\}},
    \qquad
    1+\lambda g_t(\beta,\Lambda)>0.
    \label{eq:el-kkt-weights}
\end{equation}
The remaining KKT equation is the dual moment equation
\begin{equation}
    P_T\left[\frac{g_t(\beta,\Lambda)}
    {1+\lambda g_t(\beta,\Lambda)}\right]=0.
    \label{eq:el-dual-kkt}
\end{equation}
The log empirical likelihood ratio is therefore
\begin{equation}
    \ell_T(\beta,\Lambda)
    =2\sum_{t=1}^T\log\{1+\lambda_Tg_t(\beta,\Lambda)\}.
    \label{eq:el-ratio-kkt}
\end{equation}

The quadratic EL reduction is also mostly KKT algebra. If
\(\max_t|\lambda_Tg_t|=o_p(1)\), then expanding
\eqref{eq:el-dual-kkt} gives
\[
    0=P_T\left[\frac{g_t}{1+\lambda_Tg_t}\right]
      =P_Tg_t-\lambda_TP_Tg_t^2+R_{1T},
    \qquad
    R_{1T}=o_p(T^{-1/2}).
\]
Hence, with \(V_T=P_Tg_t^2\),
\[
    \lambda_T=\frac{P_Tg_t}{V_T}+o_p(T^{-1/2}).
\]
Expanding \eqref{eq:el-ratio-kkt} then yields
\begin{equation}
    \ell_T(\beta_0,\Lambda)
    =\frac{T\{P_Tg_t(\beta_0,\Lambda)\}^2}{V_T(\Lambda)}+o_p(1)
    =\frac{S_T(\beta_0,\Lambda)^2}{V_T(\Lambda)}+o_p(1).
    \label{eq:el-kkt-quadratic}
\end{equation}
The stochastic work is not to rediscover EL; it is only to verify the small size
conditions under which the KKT expansion is valid.

\section{Population Hilbert Space}

\subsection{Population object in \texorpdfstring{\(L^2(P)\)}{L2P}}

To express breaks in population geometry, include rescaled time \(S\in(0,1]\).
Let
\[
    O=(Y,X,W,S,U)
\]
be a population random element with law \(P\). The true regime map is
\[
    j_0(S)=j
    \quad\Longleftrightarrow\quad
    \tau_{j,0}<S\le\tau_{j+1,0}.
\]
The population model is the identity in \(\LtwoP\):
\begin{equation}
    Y=\beta_0X+\mu_0(S,W)+U,
    \qquad
    \mu_0(S,W)=\sum_{j=0}^{q_0}m_j(W)
    \ind\{\tau_{j,0}<S\le\tau_{j+1,0}\}.
    \label{eq:l2p-model}
\end{equation}

\subsection{Population nuisance space}

For a population partition \(\Lambda\), define
\[
    \N_\Lambda^P
    =\overline{\left\{
      \sum_{j=0}^{q} \ind\{S\in I_j(\Lambda)\}h_j(W):
      h_j\in L^2(P_W)
    \right\}}^{L^2(P)}.
\]
Let \(\Pi_\Lambda^P\) be the \(\LtwoP\)-orthogonal projection onto
\(\N_\Lambda^P\), and let
\[
    M_\Lambda^P=I-\Pi_\Lambda^P.
\]
For \(\beta\in\mathcal B\), define the population partial outcome
\[
    Y_\beta^\circ=Y-\beta X.
\]
The population profiled residual and score direction are
\[
    r_\beta^P(\Lambda)=M_\Lambda^PY_\beta^\circ
    =M_\Lambda^P(Y-\beta X),
    \qquad
    v^P(\Lambda)=M_\Lambda^Pw.
\]
The population moment is
\begin{equation}
    \Psi(\beta,\Lambda)
    =\innerP{r_\beta^P(\Lambda)}{v^P(\Lambda)}.
    \label{eq:population-moment}
\end{equation}

\section{Population Identification}

At the true population partition \(\Lambda_0\), abbreviate
\[
    \N_0^P=\N_{\Lambda_0}^P,
    \qquad
    \Pi_0^P=\Pi_{\Lambda_0}^P,
    \qquad
    M_0^P=M_{\Lambda_0}^P,
    \qquad
    v_0^P=v^P(\Lambda_0)=M_0^Pw.
\]
Because the true nuisance is regime-specific with the true population partition,
\[
    \mu_0(S,W)\in\N_0^P.
\]
Therefore the true-partition population residual-maker removes it:
\begin{equation}
    M_0^P\mu_0=0.
    \label{eq:population-nuisance-cancellation}
\end{equation}
This is the \(L^2(P)\) analog of the finite-sample statement that the correct
break-aware nuisance space absorbs the nuisance component.

Using the population model \eqref{eq:l2p-model}, the population partial outcome
at a generic \(\beta\) is
\[
\begin{aligned}
    Y_\beta^\circ
    &=Y-\beta X \\
    &=\mu_0(S,W)+U-(\beta-\beta_0)X.
\end{aligned}
\]
Projecting onto \((\N_0^P)^{\perp_P}\) gives
\begin{equation}
\begin{aligned}
    r_\beta^P(\Lambda_0)
    &=M_0^PY_\beta^\circ \\
    &=M_0^P\mu_0+M_0^PU-(\beta-\beta_0)M_0^PX \\
    &=M_0^PU-(\beta-\beta_0)M_0^PX,
\end{aligned}
\label{eq:population-profiled-residual-identity}
\end{equation}
because \(M_0^P\mu_0=0\). Taking the inner product with the true population
score direction \(v_0^P\) gives
\begin{equation}
    \Psi(\beta,\Lambda_0)
    =\innerP{M_0^PU}{v_0^P}
     -(\beta-\beta_0)\innerP{M_0^PX}{v_0^P}.
    \label{eq:population-identification}
\end{equation}
This derivation is algebraic. The two population assumptions enter only after
\eqref{eq:population-identification} is obtained.

The exogeneity condition is
\[
    \innerP{M_0^PU}{v_0^P}=0.
\]
The relevance condition is
\[
    \Gamma=\innerP{M_0^PX}{v_0^P}\ne0.
\]
Under these two conditions,
\[
    \Psi(\beta,\Lambda_0)=-(\beta-\beta_0)\Gamma.
\]
Thus
\begin{equation}
    \Psi(\beta,\Lambda_0)=0
    \quad\Longleftrightarrow\quad
    \beta=\beta_0.
    \label{eq:unique-identification}
\end{equation}
\part{Small Stochastic Bounds}

\section{Small Stochastic Bounds}

Part I has already fixed the objects. For every candidate partition \(\Lambda\),
write
\[
    \bm r_\beta(\Lambda)=\bm M_{K,\Lambda}(\bm y-\beta\bm x),
    \qquad
    \bm a(\Lambda)=\bm M_{K,\Lambda}\bm w,
\]
and
\[
    g_t(\beta,\Lambda)=r_{\beta,t}(\Lambda)a_t(\Lambda),
    \qquad
    \Psi_T(\beta,\Lambda)=P_Tg(\beta,\Lambda)
    =\innerT{\bm r_\beta(\Lambda)}{\bm a(\Lambda)}.
\]
At the true partition, abbreviate
\[
    \bm M_0=\bm M_{K,\Lambda_0},
    \qquad
    \bm v_0=\bm M_0\bm w,
    \qquad
    \bm y_\beta^\circ=\bm y-\beta\bm x.
\]
For the estimated partition, write
\[
    \widehat{\bm M}=\bm M_{K,\widehat\Lambda},
    \qquad
    \widehat{\bm v}=\widehat{\bm M}\bm w.
\]
The variance objects are
\[
    V_T(\Lambda)=P_Tg_t(\beta_0,\Lambda)^2,
    \qquad
    V_T^0=V_T(\Lambda_0),
    \qquad
    \widehat V_T=V_T(\widehat\Lambda).
\]

The stochastic part proves the following targets one by one:
\begin{equation}
    \sup_{\beta\in\mathcal B}
    |\Psi_T(\beta,\widehat\Lambda)-\Psi(\beta,\Lambda_0)|=\op(1),
    \label{eq:uniform-bridge}
\end{equation}
\[
    \sqrt T\Psi_T(\beta_0,\widehat\Lambda)\Rightarrow N(0,V),
    \qquad
    \widehat V_T\to_p V>0,
\]
and the local empirical-likelihood conditions that make the KKT expansion
\eqref{eq:el-kkt-quadratic} valid. None of these statements changes the
geometric objects. They only control their sampling error.

\subsection{Standing stochastic conditions}

Let \(\mathcal L_T\) denote the admissible set of candidate partitions searched
by the estimator. For local break-date arguments, let
\(\mathcal N_T(r_T)\subset\mathcal L_T\) denote the event/set of partitions whose
break dates are within the localization radius \(r_T\) of \(\Lambda_0\). The
following assumptions are the stochastic input for Part II. They are stated in
the workbook notation so each proof below can be read directly.

\begin{assumption}[Empirical-to-population inner-product bridge]
\label{ass:gram-stability}
Uniformly over \(\Lambda\in\mathcal L_T\), the empirical Gram matrices of the
block sieve design are stable:
\[
    \sup_{\Lambda\in\mathcal L_T}
    \left\|T^{-1}\bm Q_{K,\Lambda}'\bm Q_{K,\Lambda}-G_Q(\Lambda)\right\|_{op}
    =\op(1),
\]
where
\[
    0<c\le \lambda_{\min}\{G_Q(\Lambda)\}
    \le \lambda_{\max}\{G_Q(\Lambda)\}\le C<\infty
    \qquad \text{uniformly over }\Lambda\in\mathcal L_T.
\]
The same uniform empirical laws hold for the projected cross moments involving
\(\bm Q_{K,\Lambda}\), \(\bm x\), \(\bm w\), and \(\bm u\). Equivalently, for
the finite sieve and break classes used here,
\[
    \sup_{f,g\in\mathcal F_{K,T}}
    |\innerT{\bm f_T}{\bm g_T}-\innerP{f}{g}|=\op(1).
\]
This is the bridge between \(L^2(P_T)\) and \(L^2(P)\). The spaces are not being
identified as the same space; only the relevant inner products converge.
\end{assumption}

\begin{assumption}[Score-negligible sieve remainder]
\label{ass:score-negligible-sieve-remainder}
For the true nuisance vector \(\bm\mu_0\), the true-partition finite-sieve
remainder
\[
    \bm r_{\mu,T}=\bm M_0\bm\mu_0
\]
satisfies
\[
    \|\bm r_{\mu,T}\|_T=\op(1),
    \qquad
    \sqrt T\,|\innerT{\bm r_{\mu,T}}{\bm v_0}|=\op(1),
    \qquad
    P_T\{r_{\mu,T,t}^2v_{0,t}^2\}=\op(1).
\]
The first bound is the consistency-scale approximation bound. The second is the
score-scale bias bound. The third removes the sieve remainder from the variance.
\end{assumption}

\begin{assumption}[Estimated-partition primitive closure]
\label{ass:directional-break-replacement}
The RSS break estimator localizes at radius \(r_T\):
\[
    \Pr\{\widehat\Lambda\in\mathcal N_T(r_T)\}\to1.
\]
On \(\mathcal N_T(r_T)\), local perturbations of the projected score and second
moment are negligible:
\[
    \sup_{\Lambda\in\mathcal N_T(r_T)}
    \sqrt T\,|\Psi_T(\beta_0,\Lambda)-\Psi_T(\beta_0,\Lambda_0)|=\op(1),
\]
\[
    \sup_{\Lambda\in\mathcal N_T(r_T)}
    |V_T(\Lambda)-V_T^0|=\op(1),
\]
and, at the unscaled consistency level,
\[
    \sup_{\Lambda\in\mathcal N_T(r_T)}\sup_{\beta\in\mathcal B}
    |\Psi_T(\beta,\Lambda)-\Psi_T(\beta,\Lambda_0)|=\op(1).
\]
A sufficient local perturbation route is
\[
    \sup_{\Lambda\in\mathcal N_T(r_T)}
    T^{-1/2}|\bm r_{\mu,T}'(\bm M_{K,\Lambda}-\bm M_0)\bm w|=\op(1),
\]
\[
    \sup_{\Lambda\in\mathcal N_T(r_T)}
    T^{-1/2}|\bm u'(\bm M_{K,\Lambda}-\bm M_0)\bm w|=\op(1),
\]
with the variance replacement below. For a one-break local window, a standard
sufficient rate is
\[
    \zeta_K\sqrt{\frac{Kr_T}{T}}\to0,
\]
where \(\zeta_K\) controls the sieve-envelope size.
\end{assumption}

\begin{assumption}[RSS concentration and penalty rates]
\label{ass:rss-localization-penalty}
Let \(\Delta_\Lambda=\bm M_{K,\Lambda}-\bm M_0\). Outside the local set
\(\mathcal N_T(r_T)\), the deterministic RSS drift dominates stochastic RSS
fluctuation:
\[
    \inf_{\Lambda\in\mathcal L_T\setminus\mathcal N_T(r_T)}
    T^{-1}\bm\mu_0'\Delta_\Lambda\bm\mu_0\ge c_T,
\]
\[
    \sup_{\Lambda\in\mathcal L_T\setminus\mathcal N_T(r_T)}
    T^{-1}|2\bm\mu_0'\Delta_\Lambda\bm u+\bm u'\Delta_\Lambda\bm u|=o_p(c_T).
\]
One primitive route for the stochastic bound is conditionally sub-Gaussian
innovations plus sieve entropy/concentration over the searched RSS class.

For unknown order, the underfitted deterministic loss and overfitted stochastic
gain satisfy
\[
    \inf_{q<q_0}\{RSS_K(q)-RSS_K(q_0)\}\ge cT\Delta_T^2+o_p(T\Delta_T^2),
\]
\[
    \sup_{q>q_0}\{RSS_K(q_0)-RSS_K(q)\}=\Op(K\log T),
\]
and the penalty obeys
\[
    K\log T=o\{\operatorname{pen}_T(q,K)\},
    \qquad
    \operatorname{pen}_T(q,K)=o(T\Delta_T^2)
\]
for every fixed nonzero order difference.
\end{assumption}

\begin{assumption}[Projected LLN and CLT at the true partition]
\label{ass:projected-lln-clt}
The true-partition projected empirical moments obey
\[
    \innerT{\bm M_0\bm u}{\bm v_0}
    -\innerP{M_0^PU}{v_0^P}=\op(1),
\]
\[
    \innerT{\bm M_0\bm x}{\bm v_0}
    -\innerP{M_0^PX}{v_0^P}=\op(1),
\]
and
\[
    T^{-1/2}\bm u'\bm v_0\Rightarrow N(0,V),
    \qquad 0<V<\infty.
\]
\end{assumption}

\begin{assumption}[Second and higher moment control]
\label{ass:moment-control}
With \(\tilde g_t^0=(\bm M_0\bm u)_t v_{0,t}\),
\[
    P_T(\tilde g_t^0)^2\to_p V,
    \qquad
    P_T|g_t(\beta_0,\Lambda_0)|^3=\Op(1),
    \qquad
    P_T|g_t(\beta_0,\widehat\Lambda)|^3=\Op(1),
\]
and
\[
    \max_{t\le T}|g_t(\beta_0,\Lambda_0)|/\sqrt T=\op(1),
    \qquad
    \max_{t\le T}|g_t(\beta_0,\widehat\Lambda)|/\sqrt T=\op(1).
\]
\end{assumption}

\begin{assumption}[Scalar convex-hull condition]
\label{ass:convex-hull}
With probability tending to one, both scalar samples
\(\{g_t(\beta_0,\Lambda_0):1\le t\le T\}\) and
\(\{g_t(\beta_0,\widehat\Lambda):1\le t\le T\}\) contain at least one strictly
positive value and at least one strictly negative value.
\end{assumption}

\subsection{Primitive stochastic assumptions sufficient for Part II}

The high-level assumptions above are implied by familiar primitive bounds. Let
\[
    v^*_{t,0}=v_{0,t}=(\bm M_0\bm w)_t,
    \qquad
    \widehat v^*_{t,0}=(\widehat{\bm M}\bm w)_t.
\]
A useful sufficient primitive package is:
\begin{enumerate}[leftmargin=2em]
\item \textbf{Envelope control.}
      \[
          \max_{t\le T}|v^*_{t,0}|=\Op(1+\zeta_K),
          \qquad
          \max_{t\le T}|\widehat v^*_{t,0}|=\Op(1+\zeta_K).
      \]
\item \textbf{Variance replacement.} If
      \(\sigma_t^2=\E(u_t^2\mid\mathcal F_{t-1})\), then
      \[
          T^{-1}\sum_{t=1}^T
          \sigma_t^2\{(\widehat v^*_{t,0})^2-(v^*_{t,0})^2\}=\op(1).
      \]
\item \textbf{Local sieve-bias stability.}
      \[
          \sup_{\Lambda\in\mathcal N_T(r_T)}
          T^{-1/2}|\bm r_{\mu,T}'(\bm M_{K,\Lambda}-\bm M_0)\bm w|=\op(1).
      \]
\item \textbf{Local noise perturbation.}
      \[
          \sup_{\Lambda\in\mathcal N_T(r_T)}
          T^{-1/2}|\bm u'(\bm M_{K,\Lambda}-\bm M_0)\bm w|=\op(1),
      \]
      for example under the rate \(\zeta_K\sqrt{Kr_T/T}\to0\) in the one-break
      local perturbation calculation.
\item \textbf{RSS concentration.} The errors have enough tail control, for
      example conditional sub-Gaussian tails, so the searched RSS class satisfies
      the concentration bounds in Assumption~\ref{ass:rss-localization-penalty}.
\item \textbf{Penalty separation.} The unknown-order penalty dominates searched
      overfit gains but is smaller than the underfit deterministic loss:
      \[
          K\log T=o\{\operatorname{pen}_T(q,K)\},
          \qquad
          \operatorname{pen}_T(q,K)=o(T\Delta_T^2).
      \]
\end{enumerate}
These primitive conditions are not new geometry. They are the stochastic closure
that makes the geometric KKT construction usable with estimated breaks.
\subsection{RSS localization and order selection}

\begin{proposition}[RSS localization]
\label{prop:part2-rss-localization}
Under Assumption~\ref{ass:rss-localization-penalty}, the known-order RSS
estimator satisfies
\[
    \Pr\{\widehat\Lambda\in\mathcal N_T(r_T)\}\to1.
\]
\end{proposition}

\begin{proof}
For any candidate \(\Lambda\), Section~7 gives the exact decomposition
\[
\begin{aligned}
    RSS_K(\beta_0,\Lambda)-RSS_K(\beta_0,\Lambda_0)
    &=\bm\mu_0'\Delta_\Lambda\bm\mu_0
      +2\bm\mu_0'\Delta_\Lambda\bm u
      +\bm u'\Delta_\Lambda\bm u,
\end{aligned}
\]
where \(\Delta_\Lambda=\bm M_{K,\Lambda}-\bm M_0\). On
\(\mathcal L_T\setminus\mathcal N_T(r_T)\), Assumption~\ref{ass:rss-localization-penalty}
gives
\[
    T^{-1}\bm\mu_0'\Delta_\Lambda\bm\mu_0\ge c_T
\]
uniformly, while
\[
    T^{-1}|2\bm\mu_0'\Delta_\Lambda\bm u+\bm u'\Delta_\Lambda\bm u|=o_p(c_T)
\]
uniformly. Therefore
\[
    \inf_{\Lambda\in\mathcal L_T\setminus\mathcal N_T(r_T)}
    \{RSS_K(\beta_0,\Lambda)-RSS_K(\beta_0,\Lambda_0)\}>0
\]
with probability tending to one. A minimizer over the known-order search set
cannot lie outside \(\mathcal N_T(r_T)\) on this event.
\end{proof}

\begin{proposition}[Unknown-order penalty selection]
\label{prop:part2-order-selection}
Under Assumption~\ref{ass:rss-localization-penalty}, the penalized RSS selector
satisfies
\[
    \Pr(\widehat q=q_0)\to1.
\]
\end{proposition}

\begin{proof}
For underfitting, Assumption~\ref{ass:rss-localization-penalty} gives
\[
    \inf_{q<q_0}\{RSS_K(q)-RSS_K(q_0)\}
    \ge cT\Delta_T^2+o_p(T\Delta_T^2).
\]
The penalty difference for a fixed lower order is
\(o(T\Delta_T^2)\), so the deterministic loss from omitting a true break
dominates any penalty saving. Thus underfitted orders are rejected with
probability tending to one.

For overfitting,
\[
    \sup_{q>q_0}\{RSS_K(q_0)-RSS_K(q)\}=\Op(K\log T).
\]
The penalty increment satisfies \(K\log T=o\{\operatorname{pen}_T(q,K)\}\), so
any spurious RSS gain is dominated by the penalty increase with probability
tending to one. Thus overfitted orders are rejected with probability tending to
one. Combining the underfit and overfit cases gives \(\Pr(\widehat q=q_0)\to1\).
\end{proof}
\subsection{Gram stability}

\begin{lemma}[Gram stability gives well-behaved projections]
\label{lem:part2-gram-stability}
Under Assumption~\ref{ass:gram-stability}, uniformly over
\(\Lambda\in\mathcal L_T\),
\[
    \lambda_{\min}\{T^{-1}\bm Q_{K,\Lambda}'\bm Q_{K,\Lambda}\}\ge c/2
\]
with probability tending to one. Hence the coordinate formula for
\(\Pi_{K,\Lambda,T}\) is stable on \(\mathcal L_T\), and the residual-maker
\(\bm M_{K,\Lambda}=I-\Pi_{K,\Lambda,T}\) remains an orthogonal contraction in
\(\LtwoT\):
\[
    \|\bm M_{K,\Lambda}\bm z\|_T\le \|\bm z\|_T
    \qquad \text{for every }\bm z\in\LtwoT.
\]
\end{lemma}

\begin{proof}
By Weyl's eigenvalue inequality,
\[
\begin{aligned}
\lambda_{\min}\{T^{-1}\bm Q_{K,\Lambda}'\bm Q_{K,\Lambda}\}
&\ge
\lambda_{\min}\{G_Q(\Lambda)\}
-\left\|T^{-1}\bm Q_{K,\Lambda}'\bm Q_{K,\Lambda}-G_Q(\Lambda)\right\|_{op}.
\end{aligned}
\]
Taking the infimum over \(\Lambda\in\mathcal L_T\), Assumption~\ref{ass:gram-stability}
puts the first term above \(c\) and the second term below \(c/2\) with
probability tending to one. This proves the eigenvalue bound.

On that event, the least-squares projection has the usual stable coordinate
form
\[
    \Pi_{K,\Lambda,T}
    =\bm Q_{K,\Lambda}
      (\bm Q_{K,\Lambda}'\bm Q_{K,\Lambda})^{-1}
      \bm Q_{K,\Lambda}'.
\]
Independently of coordinates, \(\Pi_{K,\Lambda,T}\) is the orthogonal projection
onto \(\N_{K,\Lambda,T}\). Therefore
\(\bm M_{K,\Lambda}=I-\Pi_{K,\Lambda,T}\) is also an orthogonal projection.
Orthogonal projections are contractions in Hilbert norm.
\end{proof}

\subsection{Empirical-to-population projection bridge}

\begin{lemma}[Inner-product bridge from \(L^2(P_T)\) to \(L^2(P)\)]
\label{lem:part2-inner-product-bridge}
Under Assumption~\ref{ass:gram-stability}, uniformly over the sieve and break
classes used in the projected score,
\[
    \innerT{\bm f_T}{\bm g_T}=\innerP{f}{g}+\op(1).
\]
In particular,
\[
    \sup_{\Lambda\in\mathcal L_T}
    \left\|T^{-1}\bm Q_{K,\Lambda}'\bm Q_{K,\Lambda}-G_Q(\Lambda)\right\|_{op}
    =\op(1),
\]
and the corresponding projected \((x,w,u)\) cross moments converge to their
population limits.
\end{lemma}

\begin{proof}
The first display is the uniform empirical law in
Assumption~\ref{ass:gram-stability}. Taking \(f\) and \(g\) to be two
coordinates of the block sieve row gives the displayed Gram convergence.
Taking one or both functions to be the projected \(x\), \(w\), or \(u\)
directions gives the corresponding cross-moment convergence. This is exactly
the bridge needed by the workbook: finite-sample projection geometry is written
in \(L^2(P_T)\), while identification is written in \(L^2(P)\), and the common
object is convergence of the relevant inner products.
\end{proof}
\subsection{Sieve remainder is small}

\begin{lemma}[True-partition sieve remainder is negligible]
\label{lem:part2-sieve-remainder}
Under Assumption~\ref{ass:score-negligible-sieve-remainder},
\[
    \innerT{\bm M_0\bm\mu_0}{\bm v_0}=\op(1),
    \qquad
    \sqrt T\innerT{\bm M_0\bm\mu_0}{\bm v_0}=\op(1),
\]
and
\[
    V_T^0
    =P_T\{(\bm M_0\bm u)_tv_{0,t}\}^2+\op(1).
\]
\end{lemma}

\begin{proof}
By definition \(\bm r_{\mu,T}=\bm M_0\bm\mu_0\). The first two displays are
exactly the consistency-scale and score-scale components of
Assumption~\ref{ass:score-negligible-sieve-remainder}.

At \(\beta_0\),
\[
    \bm M_0\bm y_{\beta_0}^\circ
    =\bm M_0(\bm\mu_0+\bm u)
    =\bm r_{\mu,T}+\bm M_0\bm u.
\]
Thus
\[
    g_t(\beta_0,\Lambda_0)
    =\{(\bm M_0\bm u)_t+r_{\mu,T,t}\}v_{0,t}.
\]
Let \(\tilde g_t^0=(\bm M_0\bm u)_tv_{0,t}\) and
\(d_t=r_{\mu,T,t}v_{0,t}\). Then
\[
    V_T^0-P_T(\tilde g_t^0)^2
    =2P_T\tilde g_t^0d_t+P_Td_t^2.
\]
The final term is \(\op(1)\) by Assumption~\ref{ass:score-negligible-sieve-remainder}.
The cross term is bounded by Cauchy-Schwarz:
\[
    |P_T\tilde g_t^0d_t|
    \le \{P_T(\tilde g_t^0)^2\}^{1/2}\{P_Td_t^2\}^{1/2}=\op(1),
\]
using the bounded second moment in Assumption~\ref{ass:moment-control}. Hence
\(V_T^0=P_T(\tilde g_t^0)^2+\op(1)\).
\end{proof}

\subsection{Break replacement}

\begin{lemma}[Estimated partition may replace the true partition]
\label{lem:part2-break-replacement}
Under Assumption~\ref{ass:directional-break-replacement},
\[
    \sup_{\beta\in\mathcal B}
    |\Psi_T(\beta,\widehat\Lambda)-\Psi_T(\beta,\Lambda_0)|=\op(1),
\]
\[
    \sqrt T\,
    |\Psi_T(\beta_0,\widehat\Lambda)-\Psi_T(\beta_0,\Lambda_0)|=\op(1),
    \qquad
    \widehat V_T-V_T^0=\op(1).
\]
\end{lemma}

\begin{proof}
By Assumption~\ref{ass:directional-break-replacement},
\(\Pr\{\widehat\Lambda\in\mathcal N_T(r_T)\}\to1\). On this event,
\(\widehat\Lambda\) is one of the partitions over which the three local
suprema in the same assumption are taken. Therefore
\[
    \sup_{\beta\in\mathcal B}
    |\Psi_T(\beta,\widehat\Lambda)-\Psi_T(\beta,\Lambda_0)|
    \le
    \sup_{\Lambda\in\mathcal N_T(r_T)}\sup_{\beta\in\mathcal B}
    |\Psi_T(\beta,\Lambda)-\Psi_T(\beta,\Lambda_0)|=\op(1),
\]
\[
    \sqrt T|\Psi_T(\beta_0,\widehat\Lambda)-\Psi_T(\beta_0,\Lambda_0)|
    \le
    \sup_{\Lambda\in\mathcal N_T(r_T)}
    \sqrt T|\Psi_T(\beta_0,\Lambda)-\Psi_T(\beta_0,\Lambda_0)|=\op(1),
\]
and
\[
    |\widehat V_T-V_T^0|
    \le
    \sup_{\Lambda\in\mathcal N_T(r_T)}|V_T(\Lambda)-V_T^0|=\op(1).
\]
The conclusion follows after adding the negligible probability of the complement
of the localization event.
\end{proof}

\subsection{Uniform moment bridge}

\begin{lemma}[Uniform projected moment bridge]
\label{lem:part2-uniform-bridge}
Under Assumptions~\ref{ass:score-negligible-sieve-remainder},
\ref{ass:directional-break-replacement}, and \ref{ass:projected-lln-clt},
\begin{equation}
    \sup_{\beta\in\mathcal B}
    |\Psi_T(\beta,\widehat\Lambda)-\Psi(\beta,\Lambda_0)|=\op(1).
    \tag{\ref{eq:uniform-bridge}}
\end{equation}
\end{lemma}

\begin{proof}
First prove the bridge at the true partition. Since
\(\bm y_\beta^\circ=\bm\mu_0+\bm u-(\beta-\beta_0)\bm x\),
\[
\begin{aligned}
    \Psi_T(\beta,\Lambda_0)
    &=\innerT{\bm M_0\bm y_\beta^\circ}{\bm v_0} \\
    &=\innerT{\bm M_0\bm\mu_0}{\bm v_0}
      +\innerT{\bm M_0\bm u}{\bm v_0}
      -(\beta-\beta_0)\innerT{\bm M_0\bm x}{\bm v_0}.
\end{aligned}
\]
The population moment from Part I is
\[
    \Psi(\beta,\Lambda_0)
    =\innerP{M_0^PU}{v_0^P}
     -(\beta-\beta_0)\innerP{M_0^PX}{v_0^P}.
\]
Subtracting gives
\[
\begin{aligned}
    \Psi_T(\beta,\Lambda_0)-\Psi(\beta,\Lambda_0)
    &=\innerT{\bm M_0\bm\mu_0}{\bm v_0} \\
    &\quad+
      \left\{\innerT{\bm M_0\bm u}{\bm v_0}
      -\innerP{M_0^PU}{v_0^P}\right\} \\
    &\quad-(\beta-\beta_0)
      \left\{\innerT{\bm M_0\bm x}{\bm v_0}
      -\innerP{M_0^PX}{v_0^P}\right\}.
\end{aligned}
\]
The first term is \(\op(1)\) by Lemma~\ref{lem:part2-sieve-remainder}. The next
two bracketed terms are \(\op(1)\) by Assumption~\ref{ass:projected-lln-clt}.
Because \(\mathcal B\) is fixed and bounded, the convergence is uniform in
\(\beta\in\mathcal B\):
\[
    \sup_{\beta\in\mathcal B}
    |\Psi_T(\beta,\Lambda_0)-\Psi(\beta,\Lambda_0)|=\op(1).
\]
Finally, Lemma~\ref{lem:part2-break-replacement} gives
\[
    \sup_{\beta\in\mathcal B}
    |\Psi_T(\beta,\widehat\Lambda)-\Psi_T(\beta,\Lambda_0)|=\op(1).
\]
The triangle inequality proves the stated bridge.
\end{proof}

\subsection{Score CLT}

\begin{lemma}[Projected score CLT]
\label{lem:part2-score-clt}
Under Assumptions~\ref{ass:score-negligible-sieve-remainder},
\ref{ass:directional-break-replacement}, and \ref{ass:projected-lln-clt},
\[
    \sqrt T\Psi_T(\beta_0,\widehat\Lambda)
    \Rightarrow N(0,V).
\]
\end{lemma}

\begin{proof}
At the true partition,
\[
\begin{aligned}
    \sqrt T\Psi_T(\beta_0,\Lambda_0)
    &=\sqrt T\innerT{\bm M_0\bm y_{\beta_0}^\circ}{\bm v_0} \\
    &=T^{-1/2}\bm u'\bm v_0
      +\sqrt T\innerT{\bm M_0\bm\mu_0}{\bm v_0}.
\end{aligned}
\]
The first equality is the definition of \(\Psi_T\). The second uses
\(\bm y_{\beta_0}^\circ=\bm\mu_0+\bm u\) and the symmetry/idempotence of
\(\bm M_0\), so \(\innerT{\bm M_0\bm u}{\bm v_0}=T^{-1}\bm u'\bm v_0\).
The approximation-bias term is \(\op(1)\) by
Assumption~\ref{ass:score-negligible-sieve-remainder}. The leading term
\(T^{-1/2}\bm u'\bm v_0\) converges to \(N(0,V)\) by
Assumption~\ref{ass:projected-lln-clt}. Therefore
\[
    \sqrt T\Psi_T(\beta_0,\Lambda_0)\Rightarrow N(0,V).
\]
Lemma~\ref{lem:part2-break-replacement} gives
\[
    \sqrt T\{
    \Psi_T(\beta_0,\widehat\Lambda)-\Psi_T(\beta_0,\Lambda_0)
    \}=\op(1).
\]
Slutsky's theorem gives the result for \(\widehat\Lambda\).
\end{proof}

\subsection{Variance consistency}

\begin{lemma}[Projected variance consistency]
\label{lem:part2-variance-consistency}
Under Assumptions~\ref{ass:score-negligible-sieve-remainder},
\ref{ass:directional-break-replacement}, and \ref{ass:moment-control},
\[
    \widehat V_T\to_p V>0.
\]
\end{lemma}

\begin{proof}
Lemma~\ref{lem:part2-sieve-remainder} gives
\[
    V_T^0=P_T(\tilde g_t^0)^2+\op(1),
    \qquad
    \tilde g_t^0=(\bm M_0\bm u)_tv_{0,t}.
\]
Assumption~\ref{ass:moment-control} gives \(P_T(\tilde g_t^0)^2\to_p V\).
Therefore \(V_T^0\to_p V\). Lemma~\ref{lem:part2-break-replacement} gives
\(\widehat V_T-V_T^0=\op(1)\), so \(\widehat V_T\to_p V\). The lower bound
\(V>0\) is part of Assumption~\ref{ass:projected-lln-clt}; it rules out a
degenerate residual score direction.
\end{proof}

\subsection{Max-score and convex hull}

\begin{lemma}[EL local feasibility and max-score]
\label{lem:part2-el-feasibility}
Under Assumptions~\ref{ass:moment-control} and \ref{ass:convex-hull},
\[
    \max_{t\le T}|g_t(\beta_0,\Lambda_0)|/\sqrt T=\op(1),
    \qquad
    \max_{t\le T}|g_t(\beta_0,\widehat\Lambda)|/\sqrt T=\op(1),
\]
and the origin is in the interior of the scalar convex hull for both
\(\{g_t(\beta_0,\Lambda_0):1\le t\le T\}\) and
\(\{g_t(\beta_0,\widehat\Lambda):1\le t\le T\}\) with probability tending to
one.
\end{lemma}

\begin{proof}
The two max-score bounds are displayed in Assumption~\ref{ass:moment-control}.
For scalar moments, the convex hull of the observed values is the interval from
the sample minimum to the sample maximum. The origin is in its interior exactly
when the minimum is strictly negative and the maximum is strictly positive.
Assumption~\ref{ass:convex-hull} states that this sign condition holds with
probability tending to one for both the known-partition and estimated-partition
moment arrays.
\end{proof}

\subsection{EL KKT local expansion}

\begin{lemma}[Local EL KKT expansion]
\label{lem:part2-el-local-expansion}
Under Assumptions~\ref{ass:score-negligible-sieve-remainder}, \ref{ass:directional-break-replacement}, \ref{ass:projected-lln-clt}, \ref{ass:moment-control}, and \ref{ass:convex-hull},
with \(g_t=g_t(\beta_0,\widehat\Lambda)\),
\[
    \lambda_T=\frac{P_Tg_t}{P_Tg_t^2}+\op(T^{-1/2}),
    \qquad
    \max_{t\le T}|\lambda_Tg_t|=\op(1),
\]
and
\[
    \ell_T(\beta_0,\widehat\Lambda)
    =\frac{T\{\Psi_T(\beta_0,\widehat\Lambda)\}^2}{\widehat V_T}
     +\op(1).
\]
\end{lemma}

\begin{proof}
By Lemma~\ref{lem:part2-el-feasibility}, the scalar convex-hull condition holds
with probability tending to one, so the EL dual root \(\lambda_T\) exists in the
interior region where \(1+\lambda_Tg_t>0\). By Lemma~\ref{lem:part2-score-clt},
\[
    P_Tg_t=\Psi_T(\beta_0,\widehat\Lambda)=\Op(T^{-1/2}).
\]
By Lemma~\ref{lem:part2-variance-consistency},
\[
    P_Tg_t^2=\widehat V_T\to_p V>0.
\]
Expanding the KKT equation
\[
    0=P_T\left[\frac{g_t}{1+\lambda_Tg_t}\right]
\]
around \(\lambda_T=0\) gives
\[
    0=P_Tg_t-\lambda_TP_Tg_t^2
      +O_p\{\lambda_T^2P_T|g_t|^3\}.
\]
The third-moment bound in Assumption~\ref{ass:moment-control} implies the
remainder is \(O_p(\lambda_T^2)\). Since \(P_Tg_t=\Op(T^{-1/2})\) and
\(P_Tg_t^2\to_p V>0\), the root satisfies \(\lambda_T=\Op(T^{-1/2})\) and hence
\[
    \lambda_T=\frac{P_Tg_t}{P_Tg_t^2}+\op(T^{-1/2}).
\]
Combining \(\lambda_T=\Op(T^{-1/2})\) with the max-score bound gives
\[
    \max_{t\le T}|\lambda_Tg_t|
    \le |\lambda_T|\max_{t\le T}|g_t|=\op(1).
\]
Therefore the log EL ratio can be Taylor-expanded:
\[
\begin{aligned}
    \ell_T(\beta_0,\widehat\Lambda)
    &=2\sum_{t=1}^T\log(1+\lambda_Tg_t) \\
    &=2T\lambda_TP_Tg_t-T\lambda_T^2P_Tg_t^2
      +O_p\{T|\lambda_T|^3P_T|g_t|^3\}.
\end{aligned}
\]
The final term is \(\op(1)\), because \(\lambda_T=\Op(T^{-1/2})\) and
\(P_T|g_t|^3=\Op(1)\). Substituting
\(\lambda_T=P_Tg_t/P_Tg_t^2+\op(T^{-1/2})\) yields
\[
    \ell_T(\beta_0,\widehat\Lambda)
    =\frac{T(P_Tg_t)^2}{P_Tg_t^2}+\op(1)
    =\frac{T\{\Psi_T(\beta_0,\widehat\Lambda)\}^2}{\widehat V_T}+\op(1).
\]
\end{proof}

\subsection{Known-partition Wilks}

\begin{proposition}[Known-partition profile EL Wilks]
\label{prop:part2-known-partition-wilks}
Under the true partition \(\Lambda_0\), the profile empirical likelihood ratio
satisfies
\[
    \ell_T(\beta_0,\Lambda_0)\Rightarrow\chi_1^2.
\]
\end{proposition}

\begin{proof}
The true-partition versions of Lemmas~\ref{lem:part2-score-clt} and
\ref{lem:part2-variance-consistency} give
\[
    \sqrt T\Psi_T(\beta_0,\Lambda_0)\Rightarrow N(0,V),
    \qquad
    V_T^0\to_p V>0.
\]
The max-score and convex-hull parts of Assumptions~\ref{ass:moment-control} and
\ref{ass:convex-hull} apply to \(g_t(\beta_0,\Lambda_0)\), so the KKT expansion
from Part I gives
\[
    \ell_T(\beta_0,\Lambda_0)
    =\frac{T\{\Psi_T(\beta_0,\Lambda_0)\}^2}{V_T^0}+\op(1).
\]
Therefore
\[
    \ell_T(\beta_0,\Lambda_0)
    =\left\{
      \frac{\sqrt T\Psi_T(\beta_0,\Lambda_0)}{\sqrt{V_T^0}}
     \right\}^2+
     \op(1)
    \Rightarrow\chi_1^2.
\]
\end{proof}
\subsection{Part II consequences}

The uniform bridge gives the consistency input. If \(\widehat\beta_T\) is an
approximate sample root,
\[
    \Psi_T(\widehat\beta_T,\widehat\Lambda)=\op(1),
\]
then Lemma~\ref{lem:part2-uniform-bridge} gives
\[
    \Psi(\widehat\beta_T,\Lambda_0)=\op(1).
\]
Part I gives the unique population zero, so \(\widehat\beta_T\to_p\beta_0\).

For Wilks, Lemmas~\ref{lem:part2-score-clt},
\ref{lem:part2-variance-consistency}, and \ref{lem:part2-el-local-expansion}
give
\[
    \ell_T(\beta_0,\widehat\Lambda)
    =\frac{T\{\Psi_T(\beta_0,\widehat\Lambda)\}^2}{\widehat V_T}+\op(1),
    \qquad
    \frac{\sqrt T\Psi_T(\beta_0,\widehat\Lambda)}{\sqrt{\widehat V_T}}
    \Rightarrow N(0,1).
\]
Part III only merges these two displayed statements with the geometric KKT
reduction already proved in Part I.
\part{Merging Part}

\section{Merging Part: Moment Bridge, EL KKT, and \texorpdfstring{\(\chi_1^2\)}{chi-square}}

The merge step should not say that the unscaled moment bridge is itself
equivalent to Wilks. The correct statement is more precise:

\begin{center}
\fbox{\parbox{0.9\textwidth}{
The same projected moment \(\Psi_T\) is the interface for both targets.
Unscaled convergence of \(\Psi_T\) to \(\Psi\) merges with population uniqueness
to prove consistency. Scaled convergence of \(\Psi_T\), together with the EL KKT
quadratic reduction, merges to prove \(\chi_1^2\).}}
\end{center}

\subsection{Merge A: consistency}

Part I gives the population geometry
\[
    \Psi(\beta,\Lambda_0)
    =\innerP{M_0^PY_\beta^\circ}{v_0^P}
    =-(\beta-\beta_0)\Gamma
\]
under exogeneity and relevance. Hence
\[
    \Psi(\beta,\Lambda_0)=0
    \quad\Longleftrightarrow\quad
    \beta=\beta_0.
\]
Part II gives the stochastic bridge
\[
    \sup_{\beta\in\mathcal B}
    |\Psi_T(\beta,\widehat\Lambda)-\Psi(\beta,\Lambda_0)|=\op(1).
\]
Therefore any approximate sample root satisfies
\[
    \Psi_T(\widehat\beta_T,\widehat\Lambda)=\op(1)
    \quad\Longrightarrow\quad
    \widehat\beta_T\to_p\beta_0.
\]
This is the estimation target.

\subsection{Merge B: Wilks inference}

For inference, Part I gives the EL KKT representation
\[
    \ell_T(\beta,\widehat\Lambda)
    =2\sum_{t=1}^T\log\{1+\lambda_Tg_t(\beta,\widehat\Lambda)\},
\]
where \(\lambda_T\) solves
\[
    P_T\left[\frac{g_t(\beta,\widehat\Lambda)}
    {1+\lambda_Tg_t(\beta,\widehat\Lambda)}\right]=0.
\]
Part II verifies the local size and convex-hull conditions, so the KKT equation
can be expanded at \(\beta=\beta_0\):
\[
    \ell_T(\beta_0,\widehat\Lambda)
    =\frac{T\{\Psi_T(\beta_0,\widehat\Lambda)\}^2}
           {V_T(\widehat\Lambda)}+\op(1).
\]
Part II also gives
\[
    \frac{\sqrt T\Psi_T(\beta_0,\widehat\Lambda)}
         {\sqrt{V_T(\widehat\Lambda)}}
    \Rightarrow N(0,1).
\]
Combining the last two displays gives
\begin{equation}
    \ell_T(\beta_0,\widehat\Lambda)
    \Rightarrow \chi_1^2.
    \label{eq:target-wilks}
\end{equation}

Thus the ``equivalence'' is conditional on the geometric KKT reduction. Once
EL has been reduced to the square of the standardized projected moment, Wilks is
equivalent to the standardized score CLT. It is not equivalent to the unscaled
LLN bridge used for consistency.

\begin{center}
\fbox{\parbox{0.9\textwidth}{
\textbf{Proof-body split.} Geometry owns the model, nuisance tangent space,
projection KKT, residual score, population zero, and EL primal/dual KKT. Small
stochastic bounds only verify convergence and local size. The merging part then
turns the same projected moment into either consistency or Wilks.}}
\end{center}

\section{Three-Part Logic Checklist}

\begin{center}
\begin{tabular}{>{\raggedright\arraybackslash}p{0.22\textwidth}
                >{\raggedright\arraybackslash}p{0.35\textwidth}
                >{\raggedright\arraybackslash}p{0.33\textwidth}}
\toprule
Proof part & Object & What it proves \\
\midrule
Geometry / KKT &
\(\bm y=\beta_0\bm x+\bm\mu_0+\bm u\in\LtwoT\) &
The full sample model is a Hilbert-space identity. \\
\addlinespace
Geometry / KKT &
\(\N_{K,\Lambda,T}=\operatorname{col}(Q_{K,\Lambda})\) &
The sieve basis represents the nuisance tangent plane. \\
\addlinespace
Geometry / KKT &
\(M_{K,\Lambda}=I-\Pi_{K,\Lambda,T}\) &
Profiling is the KKT normal equation for orthogonal projection. \\
\addlinespace
Geometry / KKT &
\(g_t=\{M(y-\beta x)\}_t\{Mw\}_t\) &
The profile score is one projected inner product. \\
\addlinespace
Geometry / KKT &
EL primal/dual KKT &
EL is a convex simplex problem cut by the moment hyperplane. \\
\addlinespace
Population geometry &
\(\Psi(\beta,\Lambda_0)=-(\beta-\beta_0)\Gamma\) &
Exogeneity and relevance give the unique population zero. \\
\addlinespace
Small stochastic &
\(\sup_\beta|\Psi_T-\Psi|=\op(1)\) &
The sample projected moment reaches the population moment. \\
\addlinespace
Small stochastic &
CLT, variance, max-score, convex hull &
The EL KKT expansion is valid locally and the standardized score is normal. \\
\addlinespace
Merging &
Moment bridge plus unique zero &
Consistency: \(\widehat\beta_T\to_p\beta_0\). \\
\addlinespace
Merging &
EL KKT quadratic reduction plus CLT &
Wilks: \(\ell_T(\beta_0,\widehat\Lambda)\Rightarrow\chi_1^2\). \\
\bottomrule
\end{tabular}
\end{center}

\section{Efficiency Language}

The result proved in this workbook is orthogonal profile empirical-likelihood
Wilks inference. It is efficient-score motivated because the nuisance component
is projected out before forming the score. It is not, by itself, a full
semiparametric efficiency theorem. A full efficiency claim would additionally
require an \(L^2(P)\) approximation of the residualized score to the efficient
score and an asymptotic linear representation
\[
    \sqrt T(\widehat\beta_T-\beta_0)
    =I_{\mathrm{eff}}^{-1}T^{-1/2}\sum_{t=1}^T S_{\mathrm{eff},t}+\op(1).
\]
Without those two extra ingredients, the correct claim is profile EL Wilks
inference after orthogonalization, not semiparametric efficiency.
\section{Minimal Final Statement}

The full profile-sieve EL proof is a three-part argument.

\begin{quote}
\textbf{Geometric/KKT body.} In \(\LtwoT\), the model is
\(\bm y=\beta_0\bm x+\bm\mu_0+\bm u\). A candidate partition and sieve define the
nuisance tangent plane \(\N_{K,\Lambda,T}=\operatorname{col}(Q_{K,\Lambda})\).
Profiling is orthogonal projection onto this space. The residual score is the
projected inner product
\[
    \Psi_T(\beta,\Lambda)
    =\innerT{M_{K,\Lambda}(\bm y-\beta\bm x)}{M_{K,\Lambda}\bm w}.
\]
EL is then a convex simplex problem with the moment hyperplane
\(\sum_t\pi_tg_t=0\), whose KKT system gives
\(\pi_t=1/[T\{1+\lambda g_t\}]\).

\textbf{Small stochastic body.} The stochastic lemmas verify that the empirical
projection, estimated break partition, sieve approximation, score variance,
convex hull, and max-score terms are small enough for the sample KKT objects to
approximate their population limits.

\textbf{Merging body.} For consistency, combine
\(\Psi_T\to\Psi\) with the unique population zero. For Wilks, combine the EL KKT
quadratic reduction with the standardized score CLT:
\[
    \ell_T(\beta_0,\widehat\Lambda)
    =\frac{T\Psi_T(\beta_0,\widehat\Lambda)^2}{V_T(\widehat\Lambda)}+\op(1),
    \qquad
    \frac{\sqrt T\Psi_T(\beta_0,\widehat\Lambda)}
         {\sqrt{V_T(\widehat\Lambda)}}\Rightarrow N(0,1).
\]
Therefore \(\ell_T(\beta_0,\widehat\Lambda)\Rightarrow\chi_1^2\).
\end{quote}
\end{document}